%!TEX TS-program = pdflatex
%!encoding = UTF-8 Unicode
%%=============================================================================
%%=== ISEA-TEX-Vorlage ===
%%===   Hilfsvariablen ===
%%===      jmu 05.07.2012     ===
%%=============================================================================


%%=============================================================================
%%= Einheiten
%%=============================================================================

% Die \unit-Befehle sind veraltet und sollten nicht mehr für neue Dokumente
% verwendet werden. Stattdessen sollte \SI direkt verwendet werden.

% \newcommand{\unit}[2]{\SI{#1}{#2}}
% \newcommand{\unitfrac}[3]{\SI[per-mode=fraction]{#1}{#2\per#3}}
% \newcommand{\nunitfrac}[3]{\SI[per-mode=fraction,fraction-function=\sfrac]{#1}{#2\per#3}}

%% Kürzel für mathrm: $U_\s{GS}$ statt $U_\mathrm{GS}$
%% Die sonst oft verwendete Methode $U\lind{GS}$ führt nicht selten zu Problemen.
\newcommand{\s}[1]{\mathrm{#1}}

\DeclareSIUnit{\amperehour}{Ah}
\DeclareSIUnit{\watthour}{Wh}
\DeclareSIUnit{\amperesecond}{As}
\DeclareSIUnit{\kilowatthour}{kWh}
\DeclareSIUnit{\kWp}{kW_p}
\DeclareSIUnit{\Wp}{W_p}
\DeclareSIUnit{\cent}{ct}
\DeclareSIUnit{\year}{a}
\DeclareSIUnit{\rpm}{rpm}



%%=============================================================================
%%= Mathematik
%%=============================================================================

%% Differentialoperator
\newcommand{\diff}{\mathrm{d}}							%% upright differential operator
\newcommand{\niceddt}[1]{\ensuremath{\nicefrac{{\mathrm{d}}#1}{{\mathrm{d}}t}}}	%% Ableitung nach der Zeit


\newcommand{\tsup}{\textsuperscript}
\newcommand{\tsub}{\textsubscript}
\newcommand{\mis}[3]{\ensuremath{#1_{\text{#2}}^{\text{#3}}}}
\newcommand{\mtis}[3]{\ensuremath{\text{#1}_{\text{#2}}^{\text{#3}}}}
\newcommand{\mti}[2]{\ensuremath{\text{#1}_{\text{#2}}}}
\newcommand{\ddt}[1]{\ensuremath{\frac{\textnormal{d}#1}{\textnormal{d}t}}}
\newcommand{\sddt}[1]{\ensuremath{\sfrac{\textnormal{d}#1}{\textnormal{d}t}}}
\newcommand{\fref}[1]{Fig.~\ref{#1}}
\newcommand{\eref}[1]{Eq.~\ref{#1}}
\newcommand{\tref}[1]{Tab.~\ref{#1}}

%% Formatting of vectors, matrices etc.
\newcommand{\cmplx}[1]{\underline{#1}}									%% Underline complex quantities
\newcommand{\matr}[1]{\mathbf{#1}}											%% Use bold type for matrix quantities
\newcommand{\vect}[1]{\text{\boldmath\ensuremath{#1}}}	%% Use bold and italic type for vectors

%%= Indizes
\newcommand{\lind}[1]{_{\mathrm{#1}}}				%% tiefgestellte Indizes nicht kursiv
\newcommand{\Lind}[1]{_{#1}}								%% tiefgestellte Indizes kursiv
\newcommand{\uind}[1]{^{\mathrm{#1}}}				%% hochgestellte Indizes nicht kursiv
\newcommand{\Uind}[1]{^{#1}}								%% hochgestellte Indizes kursiv

%% si-Funktion
\DeclareMathOperator{\sinc}{sinc}


%%= Aus alter ISEA-TeX-Vorlage 
\newcommand{\kmatrix}[1]{\left[\matrix{#1}\right]}
\renewcommand{\vec}{\underline}							%% Der Unterstrich wird als Vektorzeichenverwendet

\newcommand{\mb}[1]{\quad\mbox{#1}\quad}
\newcommand{\fuer}{\mb{für}}
\renewcommand{\mit}{\mb{mit}}
\newcommand{\und}{\mb{und}}

\newcommand{\mi}[2]{\ensuremath{{#1}_{\mathrm{#2}}}}     			%math mit Index
\newcommand{\me}[2]{\ensuremath{{#1}^{\mathrm{#2}}}}     			%math mit Exp
\newcommand{\mie}[3]{\ensuremath{{#1}_{\mathrm{#2}}^{\mathrm{#3}}}} 		%math m. I. u. Exp

\newcommand{\bint}[4]{\ensuremath{\int_{#1}^{#2}{#3{\mathrm{\,d}}#4}}}
\newcommand{\uint}[2]{\ensuremath{\int#1{\mathrm{\,d}}#2}}

\newcommand{\entspricht}{\mathrel{\widehat{=}}} 									%flacher




%%=============================================================================
%%= Text
%%=============================================================================

%% Abkürzung für Typewrite-Textformatierung
\newcommand{\ttt}[1]{\texttt{#1}}

%%= Backslash im Text
\newcommand{\bs}{$\backslash$}


%%=============================================================================
%%= Farben
%%=============================================================================

%%= RWTH und ISEA-Farben =
%% Textfarbe
\definecolor{textColor}{rgb}{0, 0, 0}
\newcommand{\textColor}{\color{textColor}}	

%% Offizielle Farben, Logo "RWTH", HKS 44 (100% Sättigung)
\definecolor{rwth_logo_blue}{cmyk}{1.0, 0.5, 0.0,0.0}
%\definecolor{rwth_logo_blue}{rgb}{0.659, 0.792, 0.906}	
\newcommand{\rwthLogoBlue}{\color{rwth_logo_blue}}
\newcommand{\iseaBlue}{\color{rwth_logo_blue}}


%% Offizielle Farbe, Logounterschrift "Aachen University", HKS 44 (30% Sättigung)
\definecolor{rwth_light_blue}{cmyk}{0.45, 0.14, 0.0, 0.0}
%\definecolor{rwth_light_blue}{rgb}{0.659, 0.792, 0.906}								
\newcommand{\rwthLightBlue}{\color{rwth_light_blue}}
\newcommand{\iseaLightBlue}{\color{rwth_light_blue}}

